\section{Original histories and their geometric costs}\label{app:history-details}

All comparisons in this appendix retain one original word and its
original physical groups. A history assigns an activation time
$0\le t_c\le d_c$ to each boundary of an anchored permutation of the
fine names. At a physical position one first splits this line at its
activated boundaries and then intersects its cells with the eligible
alphabet. A competitor belongs to the class $Q_h(\omega)$ when its
label is in the anchor's eligible cell at every original slot.
For a split $e=(A;B,C)$ define the logarithmic loss
\[
 \ell_e(a)=\mathbf1_{a\in A}
   \log\frac{\pi(A\cap E)}{\pi(\operatorname{side}_e(a)\cap E)}.
\]
It is zero when only one side is eligible. If $B_e(t)$ is its sum on
tail slots, the split cost is $C_e(t)=B_e(t)-(d_e-t)$. Product-law
telescoping gives $-\log\pi(Q_h)=\sum_eB_e(t_e)$.
Use complete-cell boundary tests, processing the cell on the left
before its right boundary activates. Equivalent one-sided conventions
give the same assertions when used consistently. All original slot
representatives, including their multiplicities within a cell, remain
in the history sums.

\subsection{Legal interpolation on the original time cells}

Fix the original permutation, the weak activation order, and a physical
block for every variable cut. In a block $[a,b]$ of length $L>0$ its
assigned cuts satisfy $a\le t_1\le\cdots\le t_m\le b$ and have
deadlines at least $b$. Put $t_0=a,t_{m+1}=b$ and
\[
 w_q=(t_{q+1}-t_q)/L,\qquad0\le q\le m.
\]
The vertex $v_q$ puts the first $q$ cuts at $a$ and the remaining cuts
at $b$. The weights are nonnegative, sum to one, and
$t_r=\sum_qw_q(v_q)_r$, since
$\sum_{q<r}w_q=(t_r-a)/L$. Over the physical blocks take the product
weights $w_v$. There are at most $2^n$ vertices for $n\le H-1$ cuts,
and only dimension-many permutations, activation orders and block
assignments. Every vertex is deadline-legal, although it need not be
an available arithmetic pivot.

Reference costs are affine on each physical block. Thus, writing
$\overline M_h=\sum_e\overline C_e$,
\begin{equation}
 \overline M_h(t)=\sum_vw_v\overline M_h(v),\qquad
 \min_t\overline M_h(t)=\min_v\overline M_h(v).
 \label{history:reference-interpolation}
\end{equation}
Any common offset is subtracted once on both sides. Ties created at a
vertex are computed in the chamber's weak order; the entropy telescope
makes the total independent of another consistent resolution of those
ties. Reference nonnegativity is a separate geometric input.

For the actual word and a split assigned to $[a,b]$, define its empirical
bridge, centered at the actual block total, by
\[
 \mathcal E_e(u)=\sum_{s\in[a,u)}\ell_e(\omega_s)
       -\frac{u-a}{L}\sum_{s\in[a,b)}\ell_e(\omega_s).
\]
Direct subtraction cancels all future groups and the deadline, giving
\begin{equation}
 C_e(u)=\frac{b-u}{L}C_e(a)+\frac{u-a}{L}C_e(b)-\mathcal E_e(u),
 \quad
 \sum_eC_e(t_e)=\sum_vw_v\sum_eC_e(v_e)-\sum_e\mathcal E_e(t_e).
 \label{history:actual-interpolation}
\end{equation}
Every term uses the same actual anchor, including its old lists.
For simultaneous cuts, the chain rule gives the pointwise identity
\[
 \sum_{e\ {\rm simultaneous}}\ell_e(c)
   =\log\frac{\pi(\Pi_-(c))}{\pi(\Pi_+(c))}.
\]
They therefore form one empirical bridge of the bundled partition loss.
For a synchronous full-to-singleton split in a uniform alphabet this
loss is the constant $\log H$; its bridge is purely uncolored.

Suppose the original named counts on a block satisfy
\begin{equation}
 |N_c([s,t))-\nu_c(t-s)|
 \le\vartheta(t-s)+K[1+d(s)^\zeta+d(t)^\zeta],
 \qquad d(u)=\min(u-a,b-u),\quad0<\zeta<1.
 \label{history:named-REG}
\end{equation}
Using the shorter end interval in the empirical bridge yields
\begin{equation}
 |\mathcal E_e(u)|\le C_H\vartheta d(u)+C_HK[1+d(u)^\zeta].
 \label{history:bridge-REG}
\end{equation}
For $u-a\le L/2$, subtract the reference mean in its prefix and
whole-block terms; their linear errors sum to at most
$2\vartheta(u-a)$, and their endpoint errors to
$K[1+d(u)^\zeta]+2K(u-a)/L$. The other half uses the negative suffix.
Sum over boundedly many colors and bounded loss coefficients.
Consequently \eqref{history:actual-interpolation} has the corresponding
lower error with $\sum_ed(t_e)$ and $\sum_ed(t_e)^\zeta$.
A small vertex cost cannot automatically pay these errors; its bundled
bridge must then retain an actual path constraint.

\subsection{Original lifetimes, projection, and delay}

A current original internal cell $A$ is born at $s_A$ and split at
$t_A$. If $a(A)$ is its first original name, every other name $c\in A$
has its unactivated original boundary inside $A$, so
$t_c\ge t_A$ and $d_c\ge t_c\ge t_A$. Only $a(A)$ can expire
during its life. An anchor-containing cell loses no names; a pure
$n_A$-cell has eligible size $n_A$ or $n_A-1$. Singleton leaves
contribute no delay.

Let $q_h$ be the number of nonempty eligible cells minus one, and
$q_{\rm geom}=\#\{c\ne *:t_c\le t<d_c\}$. Every active original
boundary starts exactly one nonanchor cell. The second rank counts
those whose first name is eligible, and the first those with any
eligible name. Hence
\begin{equation}
 D_h:=\int(q_h-q_{\rm geom})
 =\sum_{A\ {\rm proper pure internal}}
        [t_A-\max(s_A,d_{a(A)})]_+\ge0.
 \label{history:delay}
\end{equation}
This holds for all original permutations, zeros and ties.
For the tail-entropy potential
\[
 F_A(t)=\sum_{s\ge t,\,\omega_s\in A}
       \log\frac{\pi_s(A\cap E(s))}{\pi_s(\omega_s)}
       -\int_t^T(|A\cap E(u)|-1)_+\,du,
\]
put
\[
 \widehat F_A(t)=F_A(t)+\int_0^t
       [(n_A-1)-(|A\cap E(u)|-1)_+]\,du.
\]
The extra increment on a pure cell's life is exactly its term in
\eqref{history:delay}, and is zero for an anchor-containing cell.
The exact split telescope can thus be written
\begin{equation}
 \sum_eC_e(t_e)=F_{A_0}(t_{\rm root})
     +\sum_{A\ {\rm proper internal}}
             [\widehat F_A(t_A)-\widehat F_A(s_A)].
 \label{history:lifetime-telescope}
\end{equation}
The root's normalization is unchanged: all original names are eligible
before its first split. One may equivalently keep unmodified potentials
and the explicit delay, but must not count both.

For a parent alphabet $O\sqcup C$, merge all child-touching cells of a
fine partition $\Pi$ into one parent star cell, retaining pure-outside
cells, and intersect its cells with $C$ for the child partition.
For a law $\nu$, write
$H_\nu(\Pi)=\sum_a\nu(a)\log(\nu(\Pi(a))/\nu(a))$ for residual
entropy. If $P(a)$ is the parent name, $Z(a)$ the fine cell and
$\overline Z(a)$ its parent cell, the conditional entropy chain rule is
\begin{equation}
 H_\pi(\Pi)=H_{\pi_p}(\Pi_p)+\pi(C)H_{\pi(\cdot\mid C)}(\Pi_c)
                  -M(\Pi),\qquad
 M(\Pi)=I(P;Z\mid\overline Z)\ge0.
 \label{history:mixed-entropy}
\end{equation}
If the child occupies one cell, or no child-touching cell contains an
outside name, the contrast is identically zero. Otherwise choose an
outside name $c$ in a mixed cell and a different child-touching cell
$A'$. In their merged parent cell $Q$, the joint atom $(c,A')$ is
zero, while its product marginal is $\pi(c)\pi(A')/\pi(Q)^2$.
If every fine probability is at least $p_*$, the elementary
relative-entropy inequality $D(P\Vert Q)\ge2\|P-Q\|_{\rm TV}^2$
gives
\begin{equation}
 M(\Pi)\ge2\pi(Q)
       [\pi(c)\pi(A')/\pi(Q)^2]^2\ge2p_*^4.
 \label{history:mixed-gap}
\end{equation}
The actual per-label contrast is
\[
 D(a)=\log\frac{\pi_p(\Pi_p(P(a)))}{\pi_p(P(a))}
 +\mathbf1_{a\in C}\log\frac{\pi(\Pi_c(a))}{\pi(a)}
 -\log\frac{\pi(\Pi(a))}{\pi(a)}.
\]
Its mean is $M$, and it is bounded. It can be negative on child names:
in a mixed cell $A$ it is
$\log[\pi(Q)\pi(A\cap C)/(\pi(C)\pi(A))]$ there, and
$\log(\pi(Q)/\pi(A))$ on outside names. Original named regularity,
applied on the dimension-many constant-partition intervals, gives
\begin{equation}
 \sum_{\rm slots}D(\omega_s)
 \ge cS_{\rm mix}-C[1+\sum_{t\ {\rm mixed endpoints}}d_g(t)^\zeta].
 \label{history:mixed-actual}
\end{equation}
Here $S_{\rm mix}$ counts only times with at least two child-touching
cells and an outside name in one of them. The original intensities
have a positive lower bound. This is a fixed-word regularity estimate,
including when the history is chosen adaptively.

Eligible unresolved ranks add exactly. If there are $u$ pure outside
cells and $v$ child-touching cells, then
\[
 |O|+|C|-u-v=(|O|+1-(u+1))+(|C|-v).
\]
Let $d_h$ be the original unactivated-boundary rank and $r_h$ the
eligible unresolved rank. Then $D_{\rm orig}=\int(d_h-r_h)\ge0$:
the number of empty original cells is at most the number of ineligible
names. This is the same expired-first delay as above. If effective
costs use eligible ranks and the two actual information terms in
\eqref{history:mixed-entropy}, their exact same-word decomposition is
\[
 -\widetilde I_h=C_p^{\rm eff}+C_c^{\rm eff}
                  +D_{\rm orig}+\sum_sD(\omega_s).
\]
Any entrance/root offset is assigned once.

The projections have lawful realizations when all outside deadlines
are at most $F$ and all child deadlines at least $F$. A fine split
with two child-touching children belongs to the child; one with a
pure-outside child belongs to the formal parent; a pure-outside split
also belongs to the parent. These assignments partition the original
event multiset. The child is the original line with outside names
deleted and each new boundary time the minimum of the intervening
original times. The intervening segments are disjoint and every child
boundary is legal. Augment the formal parent by isolating each outside
name at its deadline if necessary. Its eligible partition is unchanged:
before $F$ no child has died; afterwards only the parent star is eligible.
Every augmented nonsingleton contains only unexpired names. Represent
it as a lawful anchored line by putting the star first in star cells,
and a maximal-deadline name first in pure cells. At a split put the
child containing that anchor first. Every right-child anchor is then
unexpired. Simultaneous refinements have only dimension-many resolutions.

For these lawful projections $D_p=0$, $D_c\ge0$, and
$D_{\rm gap}=D_{\rm orig}-D_c\ge0$. Before $F$ the child delay is
zero; after $F$ the full and child eligible ranks agree, whereas each
unactivated child boundary needs an entire disjoint unactivated
original segment. Thus $d_h\ge d_c$ pointwise. Rank additivity gives
\begin{equation}
 \sum t_h=\sum t_p+\sum t_c+D_{\rm gap},\quad
 \widetilde\Phi_h=\Phi_p\Phi_c e^{-D_{\rm gap}},\quad
 \widetilde I_h=I_p+I_c-D_{\rm gap}-\sum_sD(\omega_s).
 \label{history:projection}
\end{equation}
The child uses the actual selected list.

Cancel common formal-parent and augmented-parent event times, with
multiplicity. If the erased times are $e_1\le\cdots\le e_j$ and
added deadlines $a_1\le\cdots\le a_j$, refinement at every time gives
$a_l\le e_l$ and $D_{\rm gap}=\sum_l(e_l-a_l)$. Retain original
pivot labels on all unchanged projected cuts. Endpoint regularity
bounds the number of original representatives in
$[a+l,a+l+1)$ by $C(1+l)^K$. Hence for a fixed projected history
\begin{equation}
 \sum_{\rm erased fiber}e^{-\delta D_{\rm gap}}
 \le C_H\left(C\sum_{l\ge0}(1+l)^Ke^{-\delta l}\right)^H<\infty.
 \label{history:erased-fiber}
\end{equation}
Crossing physical endpoints costs only the fixed number of groups.
No retained child pivot is renumbered as a child-only representative.

\subsection{Realizing added deadlines on the actual grid}

Assume original uncolored occupancy obeys, at complete-cell endpoints
in each block,
\[
 N([s,t))\ge\lambda_0(t-s)-K[1+d(s)^\zeta+d(t)^\zeta],
 \qquad\lambda_0>0.
\]
For a physical deadline $d$, let $r(d)$ be the right boundary of its
last occupied cell, or zero if there is none. The empty suffix from
$r(d)$ to $d$ has length at most
\[
 B=\max\{1,(2K/\lambda_0)^{1/(1-\zeta)}\}+HK/\lambda_0.
\]
Indeed every whole empty block has length at most $K/\lambda_0$;
the one possible initial partial block of length $u$ satisfies
$\lambda_0u\le K(1+u^\zeta)$, giving the first term.
Replacing at most $H$ added deadline cuts by $r(d)$ changes no
activation status at any original slot. Thus the literal comparison
classes and all their weighted masses are identical for every
acceptance $F$, even if it couples old labels, types or future words.
The multipliers change by the exact factor
\[
 \Phi_h/\Phi_{h_r}=e^{-\sum(d-r(d))}\in[e^{-HB},1].
\]
Apply this only to added parent deadlines. The projection becomes
\[
 \Phi_h=\Phi_{p_r}\Phi_c e^{-D_{\rm gap}-E_{\rm round}},
 \qquad0\le E_{\rm round}\le HB.
\]
Retained original and child pivots are unchanged. The finite assignment
choices preserve \eqref{history:erased-fiber}. Projected parent rows
must allow repeated representatives: different deadlines can share
the same last occupied cell. The interpolation, signature and geometric
bounds below hold in this enlarged positive row because they use named
cuts with ties and count tuples by a product of representative counts.
No distinct-pivot theorem is invoked on that enlarged parent row.

\subsection{Branching and the terminal floating chain}

Consider a proper owner with exterior names of relative weight one and
permanent child of weight $z$, where $1<z_-\le z\le z_+$ and
$r+z\le H$. In each original group the reference intensities are
$\kappa_g$ per exterior name and $z\kappa_g$ for the child, with
$0<\kappa_-\le\kappa_g\le\kappa_+$. Assume named regularity on
all original groups used, including old ones. On the natural complete
chain $z\ge\sqrt2$ follows from $s_m>1/2$ and the terminal weight.
These are reference intensities even after an old word has been fixed.

For an original history, let
$B_h(t)=\{*\}\cup\{c:t_c>t\}$. All its names are eligible.
Ordering the names by decreasing activation time gives a lawful
star history $h^*$ with the same named cuts and multiplier; its only
nonsingleton cell is $B_h(t)$. On the same word,
\begin{equation}
 C_h-C_{h^*}=I_{h^*}-I_h=\sum_sL_h(\omega_s),
 \label{history:branch-contrast}
\end{equation}
where
\[
 L_h(a)=\mathbf1_{a\in B_h}
       \log\frac{\pi_g(B_h)}{\pi_g(a)}
        -\log\frac{\pi_g(A_h(a))}{\pi_g(a)}.
\]
Suppose the current star cell has $u$ exterior names and the pure
original internal cells have original sizes $n_l\ge2$ and eligible
sizes $q_l=n_l-\delta_l$, $\delta_l\in\{0,1\}$.
Then $d_h=u+\sum_l(n_l-1)$, and the reference contrast per unit
time equals
\[
 \kappa_g\{E_{d_h}(z)-E_u(z)-\sum_lq_l\log q_l\}.
\]
Move the $n_l-1$ nonfirst names into the star one cell at a time.
For $k=n_l-1$ and current mass $w\ge z$, the entropy gain is
$E_k(w)-q_l\log q_l$. If the first name is alive, it is at least
\[
 E_k(z)-E_k(1)\ge(z_--1)\log(1+1/z_+)=:\gamma_z>0.
\]
If expired, it is at least
$(k+1)\log(k+1)-k\log k\ge2\log2>\gamma_z$.
Thus the contrast rate is at least $c_0$ times the number of pure
original internal cells. When that number is zero the partitions
agree, so the contrast vanishes pointwise.

Let $S_{\rm pure}$ be the sum of these cell lifetimes, and $\beta_l$
the births of the positive-length components where they are present.
Subdivide at activations and physical endpoints and apply
\eqref{history:named-REG} to the bounded coefficients. Since
$d_E(t)^\zeta\le d_E(\beta_l)^\zeta+|J_l|^\zeta$ inside a
component, choosing the slope small and absorbing
$K|J_l|^\zeta$ into its linear cost gives
\begin{equation}
 C_h\ge C_{h^*}+cS_{\rm pure}
        -C[1+\sum_l d_E(\beta_l)^\zeta].
 \label{history:proper-branching}
\end{equation}
Expired-first delay has already been included through original rank.
The remaining birth-depth error cannot be absorbed into an arbitrarily
short lifetime; it needs a retained growing wall. Boundary shifts of
one fixed cell add its regularity depth error and a bounded geometric
term, preserving the same estimate on original representatives.
This is a cost comparison and asserts no order between class posteriors.

For a terminal owner all eligible names have equal reference intensity,
and the preceding strict-child gap is unavailable. Select instead the
right endpoint $\tau$ of the last positive-length interval with a
nonsingleton eligible cell. If none exists, its entropy is zero and
$I_h=-\sum t_e$. Otherwise choose two names in that last cell that
remain eligible up to $\tau$ from the left and retain one, $c$, as
a spine name. Its ancestral cell contains both for every $t<\tau$.
It therefore has eligible size $q_{\rm sp}(t)\ge2$ and $d_c\ge\tau$.
Clip activations after $\tau$ to $\tau$; this changes no eligible
entropy and pays $\sum_e(t_e-\tau)_+$ in geometric delay.

For $t<\tau$, let $B(t)$ consist of the eligible names in the spine
cell and the nonfirst names of every other original internal cell.
This set is nested. At an off-spine split it loses the right child's
first name; at a spine split keeping the left child it does the same;
at one keeping the right child it loses the old first name if alive;
a spine first-name expiration removes that name at its deadline.
Off-spine first-name expirations change nothing. Remove all but $c$
at $\tau$.
Equivalently label the refinement tree by representatives, retaining
$c$ on its ancestral branch and the original left child on off-spine
subtrees. Assign each removed representative its split time, clipped
to its own deadline if it has expired and then to $\tau$. The time
multiset before clipping is the original one. If $\delta_{\rm sp}(t)$
indicates an expired spine first name, then
\[
 |B(t)|-1=d_{\rm orig}(t)-\delta_{\rm sp}(t),\qquad
 D^*=\int_0^\tau\delta_{\rm sp}(t)\,dt+
                         \int_\tau^Td_{\rm orig}(t)\,dt\ge0
\]
is exactly the loss of the sum of activation times.

This chain has a lawful representation anchored at the original name
$0$. If $c=0$, use decreasing removal times. Otherwise let $s_0$ be
$0$'s removal time and use the line $0,c,$ followed by all other names
in decreasing removal order; give the boundary before $c$ time $s_0$.
It is legal since $s_0\le\tau\le d_c$. Call this comparator
$h^{\rm sp}$; after $s_0$ its nonsingleton cell floats away from $0$.
The exact multiplier discrepancy remains $D^*$.

If off-spine internal cells have original sizes $n_l$ and eligible
sizes $q_l=n_l$ or $n_l-1$, their $k_l=n_l-1$ nonfirst names give
$|B|=q_{\rm sp}+\sum_lk_l$. The mean contrast is
\[
 \kappa_g\{|B|\log|B|-q_{\rm sp}\log q_{\rm sp}
                                      -\sum_lq_l\log q_l\}.
\]
Absorbing a cell into a mass $w\ge2$ gains, when its first name is
alive,
$E_k(w)-(k+1)\log(k+1)=\int_1^w\log(1+k/u)\,du
\ge\log(1+1/H)$; expiration only increases the gain. Let
$S_{\rm off}$ be the off-spine lifetime. The same regularity argument
as for \eqref{history:proper-branching} gives
\begin{equation}
 C_h\ge C_{h^{\rm sp}}+D^*+cS_{\rm off}
                  -C[1+\sum_l d_E(\beta_l)^\zeta].
 \label{history:floating}
\end{equation}
A birth erased because its first name expired has depth at most its
distance from that physical deadline, which is paid by $D^*$.
Absorb these errors into a fraction of $D^*$ and a constant. Other
births remain comparison cut times.

For an erased original time mapped to a physical endpoint, its
exponentially weighted fiber is bounded. If mapped instead to an
interior retained $\tau$, named regularity gives only
\begin{equation}
 \sum_{\text{original }t\ge\tau}e^{-a(t-\tau)}
          \le C_a[1+d_E(\tau)^\zeta].
 \label{history:terminal-fiber}
\end{equation}
There are at most $H-1$ such factors. A retained $d_E(\tau)^\gamma$
wall with $\gamma>\zeta$ will pay their polynomial. No bounded
occupancy at an arbitrary interior point is presumed.

The reference cost itself pays the floating interval. This uses
nonnegativity of every complete lawful nonincreasing reference rank
history with the same entrance normalization. In a group with $r+1$
unit names its integrand is
$g_g(d)=d-\kappa_g(d+1)\log(d+1)$.
On $[s_0,\tau)$ the floating rank lies in $1\le d\le r-1$.
Increasing or decreasing it by one on that whole interval gives two
lawful reference competitors: the rank drop at $s_0$ admits the
increase, the capacity bound holds, and both end at zero at $\tau$.
Discrete concavity gives
\[
 2g_g(d)-g_g(d-1)-g_g(d+1)
    =\kappa_g\Delta^2[(d+1)\log(d+1)]\ge\kappa_*/H.
\]
The second difference is the integral of $1/x$ against the unit
triangular kernel on $[d,d+2]$. Integrating and using nonnegative
competitor costs yields
\begin{equation}
 \overline C_{h^{\rm sp}}\ge a_H(\tau-s_0),
                      \qquad a_H=\kappa_*/(2H).
 \label{history:floating-reference}
\end{equation}
Replace $c$ by $0$ in the retained set on $[s_0,\tau)$, removing
$c$ at $s_0$ instead of $0$ and keeping all other times. This gives
a lawful permanent-$0$ star $h^0$ with the same time multiset and
reference cost. On the actual word their difference is
\[
 I_{h^{\rm sp}}-I_{h^0}
   =\sum_{[s,t)\subset[s_0,\tau)}\log|B|[N_c([s,t))-N_0([s,t))].
\]
The finite subdivision and regularity therefore give
\begin{equation}
 |C_{h^{\rm sp}}-C_{h^0}|
 \le C_H\vartheta(\tau-s_0)
       +C_K[1+d_E(s_0)^\zeta+(\tau-s_0)^\zeta].
 \label{history:anchor-comparison}
\end{equation}
The global reference and depth budgets below, rather than a separate
floating-chain closure assertion, pay this difference.

Mixed time is preserved sufficiently well by these comparisons.
The augmented proper parent has no delay, so its same-cut star has
exactly the same eligible unresolved rank. The terminal comparator
has rank at least that of the original terminal partition: before
$\tau$ they are respectively $d_{\rm orig}-\delta_{\rm sp}$ and
$d_{\rm orig}-\delta_{\rm all}$, with
$\delta_{\rm sp}\le\delta_{\rm all}$; afterwards both are zero.
Iterating eligible-rank additivity implies that every comparison child
has unresolved rank at least its original child. Thus if it has at
least two cells, so did the original. If a proper parent has no pure
internal cell, its original and star partitions agree pointwise.
Writing $M_j,M_j^*$ for the original and comparison mixed indicators,
\begin{equation}
 M_j^*(t)\le M_j(t)+\mathbf1_{\{A_j(t)>0\}},\qquad
 \sum_jS_{{\rm mix},j}(h^*)\le
     \sum_jS_{{\rm mix},j}(h)+\sum_{j<m}S_{{\rm pure},j}.
 \label{history:mixed-domination}
\end{equation}
Indeed mixedness means precisely that the merged parent star has an
outside name and the child has at least two cells. Consequently a
small fixed fraction of original mixed charge can be reserved for the
global comparator while retaining half each proper pure-lifetime
charge. This is a deterministic comparison of indicators, not of the
signed per-label contrast.

\subsection{Synchronizing cuts and sharing one endpoint budget}

Here first consider a proper coarse owner, with
$g_i(d)=d-\kappa_iE_d(z)$ and $\kappa_i\ge\kappa_*>0$.
Suppose all lawful endpoint stars on the supported actual word satisfy
$C_v\ge c_0\overline C_v-h_E$, with $\overline C_v\ge0$ and
one common normalization. In a fixed chamber, synchronize its $m$
star cuts in a block at their average
$\tau=(t_1+\cdots+t_m)/m$, and let $D_i=t_m-t_1$.
The actual before/after named sets have ranks $R,q=R-m$ and remain
legal. The rank integral is unchanged because
$\int_a^b(d(t)-q)/m\,dt=\tau-a$.
Since $g_i''(d)=-\kappa_i/(d+z)\le-\kappa_*/H$, the mean excess
over this chord is at least
\[
 \frac{\kappa_*}{2H}\int_a^b(d-q)(R-d)\,dt
              \ge\frac{\kappa_*}{2H}D_i.
\]
For one cut both sides vanish. The actual contrast is supported in
$[t_1,t_m]$; named regularity on its finitely many pieces, including
the adjacent cell at a nongrid average, bounds the error by a small
linear multiple of $D_i$ and $C[1+d_i(\tau)^\zeta+D_i^\zeta]$.
Absorb the spread power into its linear charge. Combining with
\eqref{history:proper-branching}, whose births occur at these cuts,
gives
\begin{equation}
 C_h\ge C_{\rm sync}+c\sum_iD_i+c_pS_{\rm pure}
                   -C[1+\sum_id_i(\tau_i)^\zeta].
 \label{history:synchronize}
\end{equation}
No new arithmetic pivot is assigned to $\tau_i$.

The synchronized endpoint domain is a product of two-point choices,
each whole bundle placed at its block's left or right endpoint. Let
$L_{\rm multi}$ be the interpolation of their actual costs. On the
same word,
\[
 C_{\rm sync}(\tau)=L_{\rm multi}(\tau)+\sum_i\mathcal R_i(\tau_i),
\]
where $\mathcal R_i$ is minus the empirical bridge of its actual
bundled partition loss, with zero endpoints. Its difference $\Delta_i$
between endpoint costs and its remainder are independent of the
outside endpoint choices. Thus $L_{\rm multi}$ is an affine sum of
one-coordinate functions. Shift every cost by $h_E$ once and put
\[
 a_i^*=\min_{\rm outside}(C_{i,L,\rm outside}+h_E),
 \quad b_i^*=a_i^*+\Delta_i,\quad
 \ell_i(t)=(1-u_i)a_i^*+u_i b_i^*,\quad u_i=(t-a_i)/L_i.
\]
Both endpoints are nonnegative, and the same outside choice minimizes
both. Positive interpolation gives
$L_{\rm multi}+h_E\ge\ell_i(\tau_i)$ for every $i$.
For $\theta_{\rm alloc}=1/(2H)$, set
$W_i=\theta_{\rm alloc}\ell_i+\mathcal R_i$. Since there are at
most $H$ coordinates,
\begin{equation}
 C_{\rm sync}+h_E\ge\tfrac12(L_{\rm multi}+h_E)
                                      +\sum_iW_i(\tau_i).
 \label{history:shared-budget}
\end{equation}
The single endpoint budget is thereby allocated before adding walls.

If $|g_i(R)-g_i(q)|\ge\epsilon>0$, every reference endpoint at
the more expensive side exceeds $\epsilon L_i$, and
$\ell_i(t)\ge c_0\epsilon$ times the distance from the cheaper
endpoint. Equation~\eqref{history:bridge-REG} gives
$W_i(t)\ge c\operatorname{dist}(t,\text{cheaper end})-C$.
If both reference endpoints in every legal outside context exceed
$\eta L_i$, it gives $W_i(t)\ge cL_i-C$ instead.
All remaining weak coordinates retain the actual allocated wall
\begin{equation}
 \theta_{\rm alloc}\ell_i(t)+\mathcal R_i(t)
                    \ge gd_i(t)^\gamma-h_W,\qquad\gamma>\zeta.
 \label{history:allocated-wall}
\end{equation}
On native blocks the reference endpoint classification identifies
these as full-current-input to deepest-suffix edges. Old weak blocks
require their own established charge or support. At a nongrid
synchronization time, adjacent-cell regularity adds
$C[1+d_i(t)^\zeta]$, absorbed by half the strong curve; the affine
chord changes only by a bounded amount over one cell since its
endpoint difference is $O(1+L_i)$.

Combining \eqref{history:synchronize}--\eqref{history:allocated-wall}
gives positive spread and pure-lifetime penalties, linear cheaper-end
distance penalties for strong coordinates, full-length penalties for
paid weak ones, and $(g/2)d_i(\tau_i)^\gamma$ for unpaid weak ones,
minus a fixed constant. Their exponential pivot sum is bounded.
To verify multiplicity, bin spreads and endpoint distances in unit
intervals. Every cut is within its spread of the average, so original
endpoint occupancy bounds the number of tuples by a fixed polynomial
in these quantities. Exponential or stretched-exponential penalties
sum that polynomial, and only dimension-many chambers occur. Repeated
representatives in the enlarged parent space obey the same product
upper bound. The actual original comparison class is never replaced
by a synchronized class in this argument.

\subsection{Global endpoint costs and the uncontracted root}

Use the actual common-root source of Section~\ref{lower:chapter}.
Fix throughout the global source and target arguments
$0<\zeta<\gamma<1/2$, with the regularity slope and the intrinsic
wall contractions chosen before constructing the source.
Owners have products $\sigma_j$, ratios
$u_j=\sigma_{j+1}/\sigma_j$, $d_j=u_j^{-1}-1$ and
$r_j=\sigma_j/\sigma_1$. Let $T_j$ be their native lengths and
$V_j=\sum_{a\le j}T_a$. The hierarchy gives $V_j\ll T_j$.
Original groups remain separate, including short or zero ones.
The actual tuple has root $X_1=x+\epsilon$, with
\[
 x=\sum_jr_jy_j,\qquad y_j\ge0,\qquad-C\le\epsilon\le0.
\]
On its source word the proper targets are zero, the strict exits obey
$0\le\eta_j\ll\sqrt{1+T_j}$ and $\sum_{j<m}d_j\eta_j\ll1$.
Set
\[
 a_j=\eta_j,\quad b_j=\eta_{j-1}/u_{j-1}\quad(j<m),\qquad
 a_m=x/r_m,\quad b_m=\eta_{m-1}/u_{m-1}-\epsilon/r_m,
\]
with initial incoming term zero. Then $Z_j^0=a_j-b_j$ is the exact
signed owner output, and $R_0=\sum_jZ_j^0$ is the full fine residual.
In particular the bounded root error remains in a nonnegative $b_m$.
Write $K_j(v)=C_j(v)-b_j=Z_j^0-I_j(v)$.

A fine endpoint star isolates each nonanchor fine name at a lawful
original physical endpoint, keeping all unresolved names in the permanent
anchor's cell. Its atomic projections retain the same named cuts.
Let $t_j$ be owner $j$'s last positive cut and
$s_j=\mathbf1_{\{t_j\le V_j/2\}}$. The endpoint count estimates
for the displaced source boxes give the coefficient-one inequalities
\begin{align}
 K_j(v_j)&\ge s_jZ_j^0+cH_j-C,
 \label{history:local-endpoint-H}\\
 K_j(v_j)&\ge s_jZ_j^0+cJ_j-C,
 \label{history:local-endpoint-J}
\end{align}
where
\[
 H_j=s_jt_j+(1-s_j)[V_j-t_j+L_{{\rm bad},j}+S_{{\rm good},j}],
 \qquad J_j=\overline C_{v_j}+S_{{\rm good},j}.
\]
The second retains the positive reference terms of the finite-score
proof: at early endpoints use the exact prefix expression and absorb
its $O(\sqrt{1+t_j})$ centered error into a fraction of reference cost;
at late endpoints use the exact tail and pay only its zero-tail square
root. The good-group displacement remains positive. Reference cost
controls the early prefix and late bad/zero lengths, so $H_j\ll J_j+1$.
These inequalities keep $a_m$ exactly at an early terminal endpoint and
$b_m$ exactly at a late one; there is no error of size $\sqrt{V_j}$
from unrelated full groups.

There are no added cuts in a fine star's projections. The exact entropy
chain rule gives
\[
 C_0(h):=R_0-I_h(\omega_0)=\sum_jK_j(v_j)+\sum_{j,s}D_j(\omega_s).
\]
On each genuinely mixed physical group, full named central counts and
\eqref{history:mixed-gap} give $cL_g-C\sqrt{1+L_g}$; absorb half
its mean. Nonmixed contrasts vanish identically. With
\[
 J(h)=\sum_jJ_j+\sum_jS_{{\rm mix},j},\qquad
 H(h)=\sum_jH_j+\sum_jS_{{\rm mix},j},
\]
one therefore obtains $C_0(h)\ge\sum_js_jZ_j^0+cJ(h)-C$, and
also its version with $H$.

Every adverse flag change $s_j=0,s_{j+1}=1$ has a charge at its own
scale. If $t_{j+1}\ge t_j/2$, the child's early prefix contributes
at least $V_j/4$. Otherwise, on $(t_{j+1},t_j)$ an eligible parent
exterior remains in the fine anchor cell while an immediate-child
exterior has become a separate singleton. All those child names are
eligible up to $V_j\ge t_j$. This is a genuinely mixed interval of
length greater than $V_j/4$. Thus $H\ge V_j/4$, and $J\gg V_j-C$,
for every adverse change. Expanding the outputs exactly,
\[
 \sum_js_jZ_j^0=s_m(x+\epsilon)/r_m
             +\sum_{j<m}\eta_j(s_j-s_{j+1}/u_j).
\]
The equal-one terms sum to $-\sum d_j\eta_j=O(1)$. Keep every
positive $\eta_j$ from a $1$-to-$0$ change. Each $0$-to-$1$ loss
is $O(\sqrt{1+V_j})$ and is absorbed by a small fraction of its
own linear charge. We conclude
\begin{equation}
 C_0(h)\ge s_ma_m+
    \sum_{j<m:\ s_j=1,\ s_{j+1}=0}\eta_j+cJ(h)-C.
 \label{history:enhanced-endpoint}
\end{equation}
The same proof holds when the actual mixed contrast is replaced by
any fixed sufficiently small positive multiple of mixed time.

We record the relevant same-word transport calculation. In owner $j$'s
native interior bank of length $\gg T_j$, recolor $e_j$ copies of the
immediate child's first shallow leaf to the local shallow exterior.
Put $A_j=\log S_j(\text{first})$, $b_j^{\rm flip}=u_jA_{j+1}$.
The construction gives
\begin{equation}
 b_j^{\rm flip}e_j-A_je_{j-1}=y_j+\delta_j,\quad e_0=0,\quad
 0\le\delta_j<b_j^{\rm flip},\quad e_j\ll\log(2+T_j),
 \label{history:cascade}
\end{equation}
and the original fine prior ratio is
$\log[p(c_j)/p(d_j)]=b_j^{\rm flip}-A_{j+1}$.
Thus the full residual changes by
$\Delta R=\sum_{j<m}e_j(b_j^{\rm flip}-A_{j+1})$.
At a regular endpoint bank the possible signatures and per-flip
increments of the fine cost are
\[
 \begin{array}{c|c}
 \text{owner full}&0\\
 \text{owner zero, child full}&b_j^{\rm flip}\\
 \text{owner zero, child zero}&b_j^{\rm flip}-A_{j+1}.
 \end{array}
\]
Here ``full'' means the entire relevant eligible cell, with no mixed
state at those two levels. The missing signature, owner full with a
split child, is mixed. The increments follow directly by subtracting
$\Delta I$ from the stated $\Delta R$; for instance in the middle
case $\Delta I=-A_{j+1}$. This keeps the prior change.

Any discrepancy from the ideal coefficient
$b_j^{\rm flip}s_j-A_{j+1}s_{j+1}$ is paid by
$J(h)\gg\sqrt{T_j}-C$. The exhaustive cases are: a proper parent
or child state contributes an early/bad length or good-group square
root in its genuine native or old group; a mixed state contributes the
whole group; an early full owner or child has an early prefix reaching
that group; a late zero owner has a zero tail containing it; a late
zero child has a tail containing its entire native length, which is
$\gg T_j$ by the hierarchy. Every per-flip coefficient is bounded.
Since $\log(2+T_j)\le\epsilon'\sqrt{T_j}+C_{\epsilon'}$, these
exceptional banks consume an arbitrarily small fixed part of $J$.
Regular banks keep their exact increments.

The main coefficients telescope without any monotonicity of the flags:
\[
 \sum_{j<m}e_j(b_j^{\rm flip}s_j-A_{j+1}s_{j+1})
 =\sum_{j<m}s_j(y_j+\delta_j)-s_mA_me_{m-1}.
\]
Since $r_jb_j^{\rm flip}=r_{j+1}A_{j+1}$, the actual root gives
\begin{equation}
 a_m-A_me_{m-1}
 =y_m-\sum_{j<m}(r_j/r_m)\delta_j=y_m+O_H(1).
 \label{history:terminal-cancellation}
\end{equation}
Thus every fine endpoint star on the target satisfies
\begin{equation}
 C_y(h)\ge\sum_{j:s_j=1}y_j+cH(h)-C.
 \label{history:target-endpoint}
\end{equation}
The early owners can be an arbitrary subset. Each target, including the
terminal one, has coefficient one.

For later use, if a late owner's endpoint template is not full on an
original native block $g$, its charge is at least $c\sqrt{L_g}-C$.
Indeed a zero parent pays its zero tail; a proper parent pays its bad
length or good square root; a full parent with a nonfull fine class
is genuinely mixed on the entire group. For the terminal owner the
atomic and fine alphabets agree. These charges add over the finite
set of affected groups after reducing the common coefficient.

\subsection{Intrinsic chords within a global fine history}

Fix a native canonical full-current-input to deepest-suffix edge in
owner $j$ and physical block $i=[l_i,b_i]$. First consider a global
endpoint context in which the whole immediate child remains intact
on the open block. Its varying bundle is exactly the indicated owner
exteriors, with deadlines at least $b_i$; all other cuts are fixed
lawful endpoints. The global endpoint difference and centered bridge
are precisely the atomic ones:
\[
 C_0(h^+)-C_0(h^-)=K_j(v_j^+)-K_j(v_j^-)=\Delta_i,
 \qquad
 \Delta_i=(r_i-q)L_i-\log S_i\,n_i+\log(q+z)K_{i,B}.
\]
No child clock has been resampled. Let $\ell_i^{\rm intr}$ be the
repaired chord from the two actual all-context intrinsic minima,
shifted by their fixed reserve $h_{\rm intr}$.
If the child is early, $t_{j+1}\ge b_i\ge L_i$, so both global
endpoints have a charge $cL_i+h_*$ after one fixed shift. If the
child is late, the intrinsic score upper is
\[
 K_j(v_j)+h_{\rm intr}\le C(h_{\rm intr}+s_ja_j+J_j).
\]
For an early proper owner followed by a late child,
\eqref{history:enhanced-endpoint} retains exactly $a_j=\eta_j$;
for a late owner $J_j$ suffices. At the terminal owner it retains
$s_ma_m$. Hence in these cases both global endpoints dominate a
fixed positive multiple of the intrinsic endpoints. Interpolation
gives the same chord domination.

Choose a global budget share $\theta_{\rm glob}>0$. The linearly
paid case plus \eqref{history:bridge-REG} pays a strong inward wall.
In the chord case choose in advance
$0<\theta_{\rm intr}\le c_*\theta_{\rm glob}$.
The source's actual contracted wall
$\theta_{\rm intr}\ell_i^{\rm intr}+\mathcal R_i\ge gd_i^\gamma-h_W$
then implies the global allocated wall. This is why the contractions
are fixed before the source probability construction.

Terminal coordinates require the root to remain outside this allocation.
The source supplies same-root representatives with actual local
increment and
$x_i^-\le C_AX_i^0$, $y_i^-\le C_AY_i^0$, where the deterministic
scores $X_i^0,Y_i^0$ are computed at zero terminal target.
For a terminal prefix these follow by a virtual zero-root count
completion of its actual type; for the ordinary reservoir the actual
root count stays fixed and the extra left reserve is
$O_A(\log(2+V_{\rm pre}))\ll X_i^0$, while the right stays in its
fixed box. No virtual clock is used. After subtracting $a_ms_m$ from
a completed terminal endpoint, \eqref{history:enhanced-endpoint}
dominates its zero-target score. Consequently its root-separated
endpoint chord dominates a fixed multiple of the actual source
representative chord, even if the separated endpoint difference itself
is no longer $\Delta_i$.

Now allow an arbitrary fine star. Synchronize by owner within each
original physical block: $\alpha=(j,i)$ has mean time $\tau_\alpha$,
spread $D_\alpha$, and endpoint depth $d_\alpha$.
Atomic concavity, including terminal $z=1$, gives
\[
 \sum_jK_j(h_j)\ge\sum_jK_j(h_j^{\rm sync})
       +c_s\sum_\alpha D_\alpha
       -C[1+\sum_\alpha d_\alpha^\zeta].
\]
For a fine star its mixed time has the exact formula
$S_{{\rm mix},j}=(t_j-\theta_j)_+$, where $t_j$ is the owner's
largest exterior cut and $\theta_j$ is the smallest cut of any
nonanchor fine name in its entire child. All child names are eligible
until $V_j\ge t_j$. Maxima, minima and positive parts are Lipschitz,
so synchronization changes total mixed time by at most
$C_H\sum D_\alpha$. The actual mixed contrast is at least
$c_M\sum S_{{\rm mix},j}$ minus
$C[1+\sum(d_\alpha^\zeta+D_\alpha^\zeta)]$: every nontrivial
error endpoint is an original cut, within its spread of its mean.
Choose a small fixed $\epsilon_{\rm mix}>0$ relative to $c_M,c_s$.
The fine cost therefore dominates
\[
 F(\tau)=\sum_jK_j(h_j^{\rm sync})
                  +\epsilon_{\rm mix}\sum_jS_{{\rm mix},j}(h^{\rm sync})
\]
plus a positive spread charge and the stated depth error.

In each physical block order all its synchronized owner-group times.
On this ordered chamber the mixed term is affine. Use ordered-simplex
endpoint weights within that block, and products only across distinct
blocks. If $F_v$ is the endpoint value and $a=a_m$, then exactly
\begin{equation}
 F(\tau)=a\mathcal S+\widehat L+\sum_\alpha\mathcal R_\alpha,
 \quad\mathcal S=\sum_v\omega_vs_m(v),\quad
 \widehat L=\sum_v\omega_v(F_v-as_m(v)).
 \label{history:root-separated-interpolation}
\end{equation}
The enhanced endpoint argument applies with
$\epsilon_{\rm mix}S_{\rm mix}$ and yields
\[
 F_v-as_m(v)\ge
 \sum_{j:s_j=1,s_{j+1}=0}\eta_j+cJ(v)-C,
 \quad \widehat L+C\ge c\sum_v\omega_vJ(v)\ge0.
\]
The coordinate marginal remains
$\mathbb P(v_\alpha=b_i)=(\tau_\alpha-l_i)/L_i$.
The root term $a\mathcal S$ is not included in $\widehat L$.

A proper native weak coordinate can have a partially split child;
its $\mathcal R_\alpha$ is still the atomic projection bridge.
Fix only interpolation choices in other physical blocks. Let
$t_{\rm out}$ be the largest child-owner cut there, including fixed
boundary cuts. It is at most $l_i$ or at least $b_i$.
If $t_{\rm out}\ge b_i$, the child's flag is fixed. An early child
pays $cL_i$ at every current vertex; a late child leaves either the
positive owner exit for an early owner or its $J_j$ for a late one,
so the intrinsic chord comparison applies.
If $t_{\rm out}\le l_i$ and there are child cuts in this block,
a right child cut is entirely old and its reference cost is at least
$c b_i\ge cL_i$. If all child cuts are left but the parent bundle
is right, the whole block is mixed. Thus every parent-right vertex
has charge $cL_i$, and the coordinate marginal pays
$c(\tau_\alpha-l_i)$. With no child cut in this block, its flag is
fixed: an early child gives the same mixed payment and a late one
allows the intrinsic comparison. Averaging the outside choices gives
\begin{equation}
 \widehat L+C\ge c\min\{\tau_\alpha-l_i,
                         \ell_\alpha^{\rm intr}(\tau_\alpha)\}.
 \label{history:partial-child}
\end{equation}
The first branch plus regularity pays the inward curve; the second
plus the same actual $\mathcal R_\alpha$ is supplied by the contracted
source wall. The identity $\min(A,B)+E=\min(A+E,B+E)$ completes
this case without a child independence assertion.

Strong coordinates have a fixed reference drift gap and are paid
from their cheaper endpoint. Old weak and native proper-input or
wrong-name weak coordinates have two-ended reference charge $cL_i$.
These statements apply vertex by vertex using the actual marginals.
Remaining terminal coordinates use the zero-target representative
comparison above, including the actual ordinary root count. Choose
$\theta_{\rm glob}\le1/(4H^2)$ and all source contractions small
enough in advance. For every coordinate,
\[
 \theta_{\rm glob}(\widehat L+C)+\mathcal R_\alpha
                         \ge g_0d_\alpha^\gamma-C_0.
\]
There are at most $H^2$ shares. Their sum leaves at least three quarters
of $\widehat L+C$. After absorbing the earlier $d^\zeta$ errors,
\begin{equation}
 C_0(h)\ge a\mathcal S+c\sum_v\omega_vJ(v)
                      +c\sum_\alpha(d_\alpha^\gamma+D_\alpha)-C.
 \label{history:star-interpolated}
\end{equation}
Every empirical bridge has been paid once.

\subsection{Selecting an endpoint without losing the root}

Choose in each physical block an ordered-simplex vertex of maximum
weight, with a fixed tie rule. Their product $v^\#$ is lawful and
has weight at least $(H+1)^{-H}$. A chosen endpoint side has marginal
probability at least $1/(H+1)$. If it is the nearer side, its movement
is $d_\alpha$; if farther, that marginal is $d_\alpha/L_i$, so
its movement is at most $(H+1)d_\alpha$. Original unsynchronized
cuts add their spread. In particular
\begin{equation}
 \sum_v\omega_vJ(v)\ge(H+1)^{-H}J(v^\#),\qquad
 |t_c-v_c^\#|\le C_H(d_\alpha+D_\alpha).
 \label{history:positive-vertex}
\end{equation}
A coordinatewise nearest vertex need not have positive weight bounded
below; the maximum-weight choice avoids that issue.

Suppose $0\le a\le A\log(2+T)$, where $T=V_m$. Then for every
$\epsilon>0$,
\begin{equation}
 a|\mathcal S-s_m(v^\#)|
 \le\epsilon\sum_v\omega_v\overline C_m(v)
                      +\epsilon\sum_\alpha d_\alpha^\gamma+C_\epsilon.
 \label{history:flag-payment}
\end{equation}
Only the unique physical block $[l,b]$ crossing $T/2$ can vary the
flag, unless another fixed-late cut forces all flags to zero. The
coordinate marginals and controlled movement give
$|\mathcal S-s_m(v^\#)|\le C_H\sum d_\alpha/L$ in this block,
without independence: for an early representative use the union bound
for a right choice; for a late one use one chosen-right coordinate to
bound the all-left event. If $L\ge T/4$, then
$A\log(2+T)d/L\le A\log(2+4L)L^{-\gamma}d^\gamma$,
whose coefficient is arbitrarily small at large $L$; bounded $L$
costs a constant. If $L<T/4$, both endpoints lie inside
$(T/4,3T/4)$, and every contributing endpoint has last terminal cut
there. The clipped-prefix reference lower
$\overline C_m(v)\ge c_H\min(t_v,T-t_v)$ therefore pays $c_HT$.
It absorbs $A\log(2+T)$ with an arbitrarily small share and a constant.
This proves \eqref{history:flag-payment} and keeps the full root
coefficient. Applying it to \eqref{history:star-interpolated} gives
\[
 C_0(h)\ge as_m(v^\#)+cJ(v^\#)
                     +c\sum_\alpha(d_\alpha^\gamma+D_\alpha)-C.
\]

We need a slightly stronger weighted form. Put
\[
 F_\epsilon^0(h^*)=\sum_jK_j^0(h_j^*)+
                   \epsilon\sum_{j<m}S_{{\rm mix},j}(h^*),
 \quad B(h^*)=\sum_j\overline C_j(h_j^*),
 \quad\Pi(h^*)=\sum_\alpha(D_\alpha+d_\alpha^\gamma).
\]
For fixed sufficiently small $\epsilon>0$, the same proof yields
\begin{equation}
 F_\epsilon^0(h^*)\ge as_m(v^\#)
           +c[J(v^\#)+B(h^*)+\Pi(h^*)]-C.
 \label{history:weighted-star-source}
\end{equation}
Indeed mixed-time Lipschitz and atomic synchronization apply directly
to this deterministic mixed-length term. At endpoints any fixed
positive mixed coefficient pays adverse flag changes. Reserve positive
shares of reference and spread, in addition to all bridge allocations:
\[
 0\le B(h^*)-B(h^{\rm sync})\le C_H\sum_\alpha D_\alpha.
\]
The lower is concavity; the upper follows because the rank profiles
differ only between each group's first and last cut, with uniformly
bounded integrands. These shares retain $B(h^*)$, while further shares
pay $J$, curves, spread and the flag interpolation. All allocations
are fixed before constructing the source.

\subsection{From a general fine history to its source and target costs}

\begin{proposition}[Original-history source geometry]
\label{history:source}
On the source support just constructed, every lawful original fine
history $h$ has a lawful comparison star $h^*$ and endpoint
representative $v=v^\#$ such that
\begin{equation}
 C_0(h)\ge as_m(v)+c\Lambda(h)-C,
 \quad
 \Lambda=J(v)+B(h^*)+\Pi(h^*)+D_{\rm gap}+D_*+
                                   S_{\rm pure}+S_{\rm off}\ge0.
 \label{history:source-geometry}
\end{equation}
For every fixed $\delta>0$, the associated geometric envelope has a
uniformly bounded sum over all original histories and original pivot
representatives. All constants are independent of length ratios and
strict product gaps.
\end{proposition}
\begin{proof}
Recursively apply the lawful projection \eqref{history:projection},
keeping original time labels and selected lists. The exact identity is
\[
 C_0(h)=\sum_jK_j^0(h_j)+D_{\rm gap}+M_{\rm orig},
\]
with $M_{\rm orig}\ge c_MS_{\rm orig}
-C[1+\sum_{\text{original cuts }t}d_E(t)^\zeta]$ from
\eqref{history:mixed-actual}. Apply
\eqref{history:proper-branching} to proper owners and
\eqref{history:floating}, \eqref{history:anchor-comparison} to the
terminal one. Their signed output normalizations cancel in each
same-word difference. They retain positive proper pure lifetimes,
terminal off-spine lifetimes and true clipping delay $D_*$.
The remaining terminal loss is
\[
 C_H\vartheta T_{\rm float}+
 C_K[1+T_{\rm float}^\zeta+\sum_{\rm births}d_E(\beta)^\zeta],
 \qquad B(h^*)\ge a_HT_{\rm float}.
\]
If there is no terminal nonsingleton interval, choose its comparison
cuts zero, $D_*=$ the sum of its original times, and
$S_{\rm off}=T_{\rm float}=0$; this case is exact.
By \eqref{history:mixed-domination}, a sufficiently small
$\epsilon<\min(c_M/2,c_p/4)$ retains
$\epsilon S_{\rm mix}(h^*)$ while leaving positive shares of
original mixed and pure time. The weighted quantity in
\eqref{history:weighted-star-source} is therefore available.

Every remaining depth error has a specific original charge. An erased
projection event is paired to a physical deadline with its distance in
$D_{\rm gap}$. A terminal event is either retained, clipped to its
deadline, or clipped to $\tau$. In the last case
$d_E(e)\le d_E(\tau)+(e-\tau)$, with $e-\tau$ in $D_*$;
$\tau$ is a retained terminal cut when it is not a physical endpoint.
The same alternatives treat expired-first births. Retained cuts lie
within their group spread of its synchronized time. Thus all errors
are bounded by
\begin{equation}
 C[1+D_{\rm gap}^\zeta+D_*^\zeta+
                    \sum_\alpha(d_\alpha^\zeta+D_\alpha^\zeta)].
 \label{history:all-depth-errors}
\end{equation}
Choose the regularity slope small enough that a reserved fraction of
$B(h^*)\ge a_HT_{\rm float}$ pays the linear floating term.
Young's inequality pays its sublinear term with another fraction.
Use positive delays and spread for their $\zeta$-powers, and the
$\gamma>\zeta$ inward curves for $d_\alpha^\zeta$. Retain a fixed
positive share of every budget. This proves
\eqref{history:source-geometry}.

The projection fibers sum by \eqref{history:erased-fiber}. Terminal
fibers at a physical endpoint sum by the same bound; at interior
$\tau$ they incur the finite polynomial in
\eqref{history:terminal-fiber}, paid by the retained inward and
spread budgets. Remaining cuts lie in the endpoint strips
\eqref{history:positive-vertex}; their polynomial multiplicity is
summed by exponential spread and stretched-exponential depth weights.
Finite assignment chambers, names, tie resolutions and permutations
only change the constant. Generated repeated deadline representatives
remain in the enlarged positive row. This proves summability without
replacing an original class by a comparison class.
\end{proof}

\begin{proposition}[Original-history target geometry]
\label{history:target}
For the literal transported bounded density $F_y$ on the same actual
root tuple, every word of positive weight and every original fine
history satisfy, with the constructions above,
\begin{equation}
 C_y(h):=R_y-I_h(\omega_y)
       \ge\sum_{j:s_j(v)=1}y_j+c\Lambda(h)-C.
 \label{history:target-geometry}
\end{equation}
Its geometric envelope has the same original-pivot summability.
\end{proposition}
\begin{proof}
Fix any retained source preimage of the target word. Such a preimage
exists for every positive value of the literal transported density.
For a comparison fine star define
$F_\epsilon^q=\sum_jK_j^q+\epsilon S_{\rm mix}$ for $q=0,y$.
Mixed time depends only on names and cuts, so exactly
\[
 \Delta F_\epsilon(h^*)=\Delta R-\sum_j\Delta I_j(h_j^*).
\]
Each atomic information is a sum over original slots, zero when a
slot does not route to that owner; the formula remains exact when
recoloring changes child membership. Compare this change with the
change at the endpoint representative $v$ on the same source/image
pair. The residual $\Delta R$ cancels. The bounded atomic mark
vector changes only if a flipped slot meets a movement strip.
Bank $j$ has depth $\gg T_j$ in its original block; such an
intersection forces $\Pi(h^*)\ge c_HT_j^\gamma-C_H$.
Since it contains $O(\log(2+T_j))$ flips, for every fixed $\alpha>0$
\begin{equation}
 |\Delta F_\epsilon(h^*)-\Delta F_\epsilon(v)|
                      \le\alpha\Pi(h^*)+C_\alpha.
 \label{history:strip-transfer}
\end{equation}

At a regular endpoint bank, the three per-flip increments computed
above apply to $\sum_jK_j$ itself: the relevant mixed contrasts vanish
identically, ancestor exteriors have expired, and deeper owners read
neither changed shallow name. Exceptional signatures have
$J(v)\gg\sqrt{T_j}-C$, by the exhaustive endpoint cases already
proved. Their logarithmic number consumes any small fixed fraction
of $J(v)$. Hence
\[
 \Delta F_\epsilon(v)\ge
 \sum_{j<m}e_j(b_j^{\rm flip}s_j-A_{j+1}s_{j+1})
                                  -\alpha J(v)-C_\alpha.
\]
Apply \eqref{history:weighted-star-source},
\eqref{history:strip-transfer}, and the exact telescope
\eqref{history:terminal-cancellation}; nonnegative $\delta_j$ only
help. For small $\alpha$,
\begin{equation}
 F_\epsilon^y(h^*)\ge\sum_{j:s_j(v)=1}y_j+
                          c[J(v)+B(h^*)+\Pi(h^*)]-C.
 \label{history:weighted-star-target}
\end{equation}
The terminal root has coefficient one, without a comparison of different
root tuples or a monotonicity claim for target local minima.

The target retains the original named regularity with prescribed spare
margins. An interval meeting a flipped interior bank either has length
$\gg T_j$ or an endpoint of depth $\gg T_j$; otherwise it cannot
reach that bank. A spare linear slope in the first case, or a spare
$\zeta$-depth coefficient in the second, pays
$O(\log(2+T_j))$ flips after fixing a sufficiently large owner
threshold. Routed atomic counts are sums of finitely many fine counts.
Thus choose target regularity constants first and construct the source
inside them; no clock or reference intensity changes.

Project the original history on this target word. Its exact identity
is $C_y(h)=\sum_jK_j^y(h_j)+D_{\rm gap}+M_y$.
Use the proper and terminal same-word comparisons and mixed-time
domination exactly as in the preceding proof. The available weighted
budget is now \eqref{history:weighted-star-target}. The floating
reference cost pays the linear and sublinear floating errors;
\eqref{history:all-depth-errors} is paid by the same true delays,
spread and inward curves. Positive shares remain, proving
\eqref{history:target-geometry}. Fiber counting is unchanged and
keeps every original representative. This argument holds simultaneously
for all histories on the chosen source/image pair; it does not compare
class posteriors.
\end{proof}

\subsection{The actual comparison class and the complete row}

Set, for the predeclared environmental costs and count widths,
\[
 \begin{gathered}
 D_j^{\rm own}=w_j\prod_{i:\operatorname{owner}(i)=j}D_i\quad(j<m),
 \qquad D_m^{\rm own}=\prod_{i:\operatorname{owner}(i)=m}D_i,\\
 y_j=-\log D_j^{\rm own},\qquad
 D_{\rm all}=\prod_jD_j^{\rm own}.
 \end{gathered}
\]
These targets are fixed before reading the private word. Fix its actual
class $Q_h(\omega)$ and the bounded transported acceptance $F=F_y$.
It has the actual root, named count support, strict exits, reconstructed
source tests, and bounded full relative likelihood established in the
transport construction.

In a physical block $g=[l_g,b_g]$, choose an integer endpoint-strip
radius $r_g\ge1$ dominating, by a dimension factor, the depths of
all original cuts and all movement strips from $h^*$ to $v$; cap it
at $L_g/2$ when necessary. From the exact depth bookkeeping,
\begin{equation}
 \sum_gr_g^\gamma\le C_H(1+\Lambda).
 \label{history:strip-budget}
\end{equation}
If $r_g<L_g/4$, the original partition is constant on the middle
interval and $h^*$ there agrees with $v$. These are partition facts,
independent of label probabilities. Proper projection and star
comparison preserve eligible unresolved rank. The only possible
increase for $h^*$ is from expired terminal off-spine first names,
whose number is at most the original terminal off-spine internal-cell
count. Thus, pointwise,
\[
 0\le r(h^*)-r(h)\le m_{\rm off}.
\]
If $v$ is full but the original middle is not, this rank difference
is at least one over an interval of length $L_g/2$, so
$S_{\rm off}\ge L_g/2$. If $v$ is nonfull on a late owner's native
block, its endpoint charge is $J(v)\ge c\sqrt{L_g}-C$, as proved
above. If the strips occupy a quarter block, their own
$L_g^\gamma$ charge pays. Every late block is therefore geometrically
paid or the original class is full outside its endpoint strips.
Any missing polynomial floor on a paid block is absorbed by an
arbitrarily small fixed fraction of $\Lambda$ at its own length scale.

For a proper late owner's count bank, either those strips reach its
interior or the original class is full on the bank. In the first case
$\Lambda$ has a positive power of $T_j$ and pays the width floor
$w_j\gg(2+T_j)^{-1/2}$. Let $J_{\rm free}$ be any set of remaining
banks. Condition on the labels read by the actual class, revealing
their full labels if needed, then on pair-coarsened data in these banks
and every label outside the relevant two-name pairs. The remaining
counts are independent original binomials, on at least $cT_j$ trials
on accepting support. Their exact diagonal formulas are
$\eta_j=h_j+b_j^{\rm flip}K_j$.
All $h_j$ and target-band centers are measurable in this one fixed
sigma-field, independent of every remaining $K_j$. Uniform binomial
maximum atoms at arbitrary centers therefore give
\begin{equation}
 \mathbb P(\text{all remaining exit restrictions}\mid\text{fixed data})
                       \le C\prod_{j\in J_{\rm free}}w_j.
 \label{history:free-count-banks}
\end{equation}
Too-small pair totals have empty accepting support. Releasing the
fixed data preserves the bound. Early-owner blocks and arbitrary
observed signatures need no full-block reveal penalty.

For an unpaid selected path block, condition first on every complete
target type and all revealed strip labels. The remaining labels have
the original uniform multiset law. The reconstructed-source estimates,
uniform over every permitted inverse transfer, give the environmental
envelope
\begin{equation}
 \mathbb P(\text{necessary tests}\mid\text{types, revealed strips})
                           \le Cq_{i,r_i}(G_i).
 \label{history:free-path-banks}
\end{equation}
This is the joint two-bank density comparison, not a product of two
separately conditioned full bridges. These predeclared envelopes are
uniform over the types subsequently summed. Below their logarithmic
cutoff, strip occupancy leaves at least half the block for that joint
comparison. Above the cutoff take $q=1$ and pay the predetermined
polynomial floor of $D_i$ by the depth charge; supported bounded
blocks use $D_i=1$.

After complete types are fixed, different blocks' remaining multiset
permutations are independent. Drop all other restrictions, multiply
\eqref{history:free-path-banks}, then release the whole coupled type
measure and apply \eqref{history:free-count-banks}. This proves for
the original class
\begin{equation}
 \mathbb E_\pi[F\mid Q_h(\omega)]
 \le C\prod_{j\in J_{\rm free}}w_j
                 \prod_{i\in I_{\rm free}}q_{i,r_i}(G_i).
 \label{history:actual-class}
\end{equation}
Every type is summed before any fractional power. Only necessary tests
of $F$ have been used; no arbitrary child weight has been multiplied
by an averaged path probability.

The same environmental envelopes can be chosen for any sufficiently
small fixed $b_0>0$ so that
$q_{i,r}\le CD_i e^{b_0r^\gamma}$. Their stretched-exponential
coefficient sum is finite, so clipping and normalization into $[0,1]$
changes only a fixed constant. Their prescribed polynomial floors and
source moments persist. Choose $b_0$ small relative to the coefficient
in \eqref{history:target-geometry} and
\eqref{history:strip-budget}. Pay all missing late widths and path
floors with the local payments above. Then
\[
 \mathbb E_\pi[F\mid Q_h(\omega)]
 \le C e^{(c/2)\Lambda}\prod_{j:s_j(v)=0}D_j^{\rm own}.
\]
These choices are noncircular: the floor fixes the logarithmic target
range independently of $b_0$; that range fixes the finite geometry;
then choose $b_0$, and only afterwards construct envelopes and source.
Combining with \eqref{history:target-geometry} gives
\[
 e^{-C_y(h)}\mathbb E_\pi[F\mid Q_h(\omega)]
                      \le CD_{\rm all}e^{-c'\Lambda}.
\]
For retained compatible competitors the full likelihood ratio is
$e^{E_y(\omega)-E_y(\omega')}$, whose amplitude is bounded by the
strict-exit weighted likelihood identity. Therefore the true relative
posterior
\[
 q_h^{\rm rel}=\mathbb E_\pi
       [F(\omega')\pi(\omega)/\pi(\omega')\mid Q_h(\omega)]
\]
satisfies the same bound with a fixed extra factor. Raising to a fixed
$\delta>0$ and using original-pivot summability proves
\begin{equation}
 \sum_{h\ {\rm original}}e^{-\delta C_y(h)}(q_h^{\rm rel})^\delta
                              \le C_\delta D_{\rm all}^{\delta}.
 \label{history:full-row}
\end{equation}
The original class $Q_h$ has been retained throughout; its accepted
mass has never been replaced by that of $h^*$ or $v$.

\subsection{Bounded family weights and the arithmetic first mass}

Fix an original cell occupancy $g$, its finite legal label space
$\Omega$, and $0<\delta\le1$. An acceptance $F:\Omega\to[0,1]$
may couple all slots, but is independent of actual prime values and
invariant under within-cell slot permutations. Set
\[
 Z_F=\sum_YF(Y),\qquad
 M_h(F;Y)=\sum_{Z\sim_hY}F(Z),\qquad
 R_F(Y)=\sum_h\Phi_g(h)^\delta M_h(F;Y)^\delta.
\]
The finite weighted class identity is exactly
\begin{equation}
 \sum_h\Phi_g(h)^\delta\Psi_{h,1+\delta}(F)
                =\sum_YF(Y)R_F(Y).
 \label{history:weighted-class-identity}
\end{equation}
It is simply the expansion of the class power sum, so a pointwise
bound $R_F\le C$ on $F>0$ gives a total at most $CZ_F$, with no
second acceptance mass. For any positive product law $\pi$,
\[
 M_h(F;Y)=\frac{\pi(Q_h(Y))}{\pi(Y)}q_h^{\rm rel}(F;Y).
\]
This follows by inserting $\pi(Z)/\pi(Z)$ in the finite class sum;
bounded $F$ alone does not remove its relative likelihood. A pointwise
row for a support indicator also holds for every smaller bounded
weight, by positivity; a merely averaged row does not imply this.

Suppose the correlated prime-cell kernel applies and the retained
occupancy has squarefree supply $cW_g\le S_g\le W_g$, with all tuples
inside the required size cutoff. On harmonic tuple measure times
Lebesgue space define
$f(\mathbf a,t)=\sum_YF(Y)\mathbf1_{B_Y(\mathbf a)}(t)$,
where the original boxes have side $b=\log2$. Then
$\int f=b^kS_gZ_F\gg W_gZ_F$.
At a point of $B_Y$, every other containing box overlaps $B_Y$, so
its divisor log vector is within $b$ in each coordinate. The correlated
kernel and \eqref{history:weighted-class-identity} consequently give
$\int f^{1+\delta}\ll W_gZ_F$. H\"older on this same unnormalized
measure yields
\[
 \operatorname{meas}\{f>0\}
 \ge\frac{(\int f)^{(1+\delta)/\delta}}
          {(\int f^{1+\delta})^{1/\delta}}
 \gg W_gZ_F.
\]
Its support is in the original divisor union. Occupancies partition
the tuple space and can be summed. Fractional bounded acceptance needs
no auxiliary coin or additional thinning loss.

Finally, extra nominal data may be removed in a linear first-mass
identity when their support is already cell-measurable. If the joint
law of $(Y,\Xi)$ and a bounded acceptance $a(Y,\Xi)$ are equivariant
under within-cell permutations, and accepted realizations have an
invariant cell-label support $S(Y)$, then
$F(Y)=\mathbb E[a(Y,\Xi)\mid Y,g]$ is invariant and
$0\le F\le\mathbf1_S$. For every positive invariant reference law
$\pi_0$ and nonnegative $Y$-measurable weight $w$,
\[
 \sum_Y\pi_0(Y)w(Y)F(Y)=\mathbb E_{\rm joint}[w(Y)a(Y,\Xi)].
\]
Use $w=v/\pi_0$ for a desired counting weight $v$, and Tonelli for
an outer occupancy integral. The conditional law of $\Xi$ may depend
on $Y$. This preserves linear likelihood factors; it does not commute
conditional expectation with a fractional power. Uniform deletion of
specified-color slot subsets and count-dependent target laws are
permutation equivariant, so the transport has this property whenever
its source and boundary support do. The complete-cell construction in
the main proof retains its nested fractional quantities before this
linear arithmetic projection.
